\scanpage{63}

% manuscript line(s) 1–2
Ainsi on trouve équivalence $\underset{\substack{\wr\\3)}}{\unsure{3')|3)|1)}}\simeq2)$ d'où \deleted{\unsure{la}} (via $3)\longrightarrow1)$) la flèche $2)\longrightarrow1)$.
% manuscript line(s) 3
\figure{\begin{tikzpicture}[x=1cm,y=1cm]\draw (1,.14) .. controls (1.3,-.03) and (2.1,.24) .. (2.8,.23);\draw (0,0).. controls (1.1,-.14) and (2.6,.1)..(3.1,.07);\draw (.4,-.29).. controls (.55,-.12) and (2.42,-.11)..(2.92,-.12);\draw (1.9,.07)--(2.12,-.12);\end{tikzpicture}}
% manuscript line(s) 4–6
% slanted
% manuscript line(s) 4–6
\margin{Produit contracté de rev. quadratiques}

% manuscript line(s) 4–6
On veut maintenant définir le foncteur ``\uline{produit contracté}'' sur la catégorie des rev. quadratiques. On veut
% manuscript line(s) 7–11
\[
\left\{\begin{aligned}\mathbb L(X_1\wedge X_2)&\simeq\mathbb L(X_1)\otimes\mathbb L(X_2)\\ &\hspace{-5em}\substack{\text{(module déterminantiel}\\\text{du rev. quadratique,}\\\text{noté }\underline L\text{ ci-dessus)}}\\ \underset{\substack{\in\\\Gamma\mathbb L(X_1\wedge X_2)}}{\delta_{X_1\wedge X_2}}&=\delta_{X_1}\cdot\delta_{X_2}\end{aligned}\right.
\]

% manuscript line(s) 12–16
De plus, on veut que tout scindage de $X_1$, \deleted{\unsure{qui}} défini par $x_1\in E_1$ avec $\psi_{X_1}(x_1)=1$, et tout scindage de $X_2$, défini par $x_2\in E_2$ avec $\psi_{X_2}(x_2)=1$, en définisse un de $X_1\wedge X_2$, i.e. on veut une application
% manuscript line(s) 17
$(x_1,x_2)\longrightarrow x_1\wedge x_2$ :
% manuscript line(s) 17–18
\diagram{\begin{tikzcd}[ampersand replacement=\&] P_{X_1}\&{}\times{}\&P_{X_2}\arrow[r]\&P_{X_1\wedge X_2}\\ X_1\arrow[u,hook]\&{}\times{}\&X_2\arrow[u,hook]\arrow[r,dash,"\text{qui induise}"']\&X_1\wedge X_2\arrow[u,hook]\end{tikzcd}}

% manuscript line(s) 19–21
Si $X_1,X_2$ munis des scindages $x_1,x_2$, et si les modules \inserted{inv.} \unsure{discriminants|discriminant} sont $L_1,L_2$, alors les structures sont définies par
% manuscript line(s) 22–23
\[
\begin{aligned}b_1&\in\Gamma L_1&b_2&\in\Gamma L_1^{\otimes2}&b_1&=2x_1-T_1\\ \unsure{b_2|b_1|c_1}&\in\Gamma L_2&c_2&\in\Gamma L_2^{\otimes2}&b_2&=2x_2-T_2\end{aligned}
\]
% archive sheet number: 62
